\section{ProtCOVER}\label{sec:method}

ProtCOVER has four components: (i) family-level calibration; (ii) adaptation to descriptor shift; (iii) a support gate; and (iv) class-conditional versions of each. Its default (ii) is conditional calibration with features $(1,\hat w)$, for an estimated density ratio $\hat w$. Two variants replace it: \emph{ProtCOVER-W} reweights the families by $\hat w$, and \emph{ProtCOVER-D} conditions on the descriptors instead of $\hat w$.

\paragraph{(i) Family-level calibration.}
Let $\hat F_k=N_k^{-1}\sum_{i}\delta_{S_{k,i}}$ be the empirical distribution of calibration scores of family $k$. Given nonnegative family weights $\omega_1,\dots,\omega_{K+1}$, set $p_k=\omega_k/\sum_{j=1}^{K+1}\omega_j$ and
\begin{equation}\label{eq:threshold}
\hat\mu=\sum_{k=1}^{K}p_k\hat F_k+p_{K+1}\,\delta_{+\infty},\qquad
\hat q=\Qt_{1-\alpha}(\hat\mu),\qquad
C(x)=\{y\in\{0,1\}:s(x,y)\le\hat q\},
\end{equation}
where $\Qt_\tau(\mu)=\inf\{u:\mu((-\infty,u])\ge\tau\}$. With $\omega\equiv1$ every family, calibration or test, has the same total weight, spread evenly over its variants. The test family's weight sits at $+\infty$ because its labels, and hence its scores, are unknown: the guarantee holds for the threshold computed with its true scores, and moving them to $+\infty$ can only raise the threshold. This is the hierarchical conformal prediction (HCP) threshold of~\cite{lee2026hierarchical}, which equals the CDF pooling of~\cite{dunn2023distribution} up to the level adjustment. Variant-level split conformal instead gives every calibration variant weight one, so families with hundreds of variants set the quantile.

\paragraph{(ii) Adapting to descriptor shift: conditional calibration on an estimated density ratio.}
Equation~\eqref{eq:threshold} uses one threshold for every test family. ProtCOVER lets the threshold depend on the test family's label-free descriptor through a feature map $\Phi:\mathcal{D}\to\mathbb{R}^p$ with $\Phi_1\equiv1$, spanning $\mathcal{F}=\{\Phi(\cdot)^\top b:b\in\mathbb{R}^p\}$. With $\tau=1-\alpha$ and the pinball loss $\rho_\tau(u)=\max\{\tau u,(\tau-1)u\}$, let
\begin{equation}\label{eq:condcal}
L(b)=\sum_{k=1}^{K}\frac{1}{N_k}\sum_{i=1}^{N_k}\rho_\tau\big(S_{k,i}-\Phi(D_k)^\top b\big)-\tau\,\Phi(D_{K+1})^\top b,\qquad
\hat q=\sup\{\Phi(D_{K+1})^\top b:\ b\in\arg\min L\},
\end{equation}
with $\hat q=+\infty$ if $L$ is unbounded below or if the fitted value $\Phi(D_{K+1})^\top b$ is unbounded over $\arg\min L$, and $C(x)=\{y:s(x,y)\le\hat q\}$. This is a family-weighted quantile regression in which, as in~\eqref{eq:threshold}, each calibration family has total weight one and the test family's scores sit at $+\infty$ (the last term); with $\Phi\equiv1$ it is the HCP threshold, except that when the pooled distribution $\hat\mu$ reaches $1-\alpha$ exactly at that threshold it takes the next pooled calibration score, or $+\infty$ if there is none. $L$ is bounded below iff $-\tau\Phi(D_{K+1})=\sum_ja_j\Phi_j$ for some weights $(\tau-1)/N_{k(j)}\le a_j\le\tau/N_{k(j)}$ on the calibration scores $j$; otherwise the query's features cannot be balanced by the calibration families, for example when its estimated ratio far exceeds theirs, and the set is $\{0,1\}$: a built-in abstention. The default feature map is $\Phi=(1,\hat w)$, where $\hat w$ is an estimate of the density ratio $w=dQ/dP$ of descriptors, so the threshold can rise for families that are rare among the calibration families and frequent among the target families. Of the two variants, \emph{ProtCOVER-W} instead reweights the families in~\eqref{eq:threshold} with $\omega_k=\hat w(D_k)$, $k\le K+1$ (computable because $D_{K+1}$ is label-free), and \emph{ProtCOVER-D} uses the standardized support descriptors themselves as features and needs no target families.

The ratio comes from a logistic classifier separating calibration families (class $0$) from target families (class $1$) on standardized descriptors: $\hat w(d)=\frac{n_0}{n_1}\exp\{\hat f(d)\}$, with $\hat f$ the fitted log-odds and $n_0$, $n_1$ the class sizes, $L_2$ penalty chosen by cross-validated log loss on this membership task (no variant labels), rescaled to mean one over the fitting calibration families and clipped to $[1/c,c]$ ($c=20$). The descriptors measure sequence similarity (maximum identity of the family to any training protein over significant MMseqs2~\cite{steinegger2017mmseqs2} hits), alignment depth ($\log N_{\mathrm{eff}}$ and $N_{\mathrm{eff}}$ per position, with $N_{\mathrm{eff}}$ the effective number of sequences in the alignment) and embedding distance (cosine distance of mean-pooled ESM2~\cite{lin2023evolutionary} embeddings to the nearest training protein), together with $\log$ length and $\log$ number of scored variants. \emph{Cross-fitting:} the calibration families are split once at random into a third that fits $\hat w$ and two thirds used for calibration, and the target families into five folds; the ratio for queries in fold $j$ is fitted against the target families outside fold $j$, so each fold has its own ratio and its own guarantee, conditional on that ratio (Lemma~\ref{lem:crossfit}).

\emph{Computation.} We solve~\eqref{eq:condcal} through its dual linear program, which has one bounded variable per calibration score but only $p$ equality constraints; test families change only the right-hand side, so HiGHS~\cite{huangfu2018parallelizing} re-solves each from the previous basis (a few milliseconds). To obtain the supremum in~\eqref{eq:condcal} when $\arg\min L$ is not a single point, the test term's weight is raised from $\tau$ to $\tau+10^{-4}$: every solution of the perturbed program has a fitted value at least the supremum, and an infeasible perturbed dual (an unbounded perturbed program) gives $+\infty$ (Appendix~\ref{app:proofs}). The computed sets therefore contain those of~\eqref{eq:condcal}, and keep every guarantee by Lemma~\ref{lem:enlarge}; an independent primal computation of \auditN{} thresholds under tighter solver tolerances changes no audited set (Appendix~\ref{app:protocol}).

\paragraph{(iii) Support gate.}
The gate returns $\{0,1\}$ for every variant of a query family (it abstains) when few calibration families resemble it, measured by a kernel effective sample size (ESS) of the support descriptors below $m$, or when the ratio puts more than $\pi_{\max}$ of ProtCOVER-W's quantile mass on the query itself; its parameters are tuned on an inner split of the training families (Appendix~\ref{app:protocol}).

\paragraph{(iv) Class-conditional versions.}
Every component has a Mondrian version: for each candidate label $y$, the threshold $\hat q_y$ is computed from label-$y$ calibration variants only (weights $1/N_{k,y}$; families with $N_{k,y}=0$ dropped), and $y\in C(x)$ iff $s(x,y)\le\hat q_y$. These target $\Cov_0$ and $\Cov_1$ separately and are the versions we compare throughout.
